\documentclass[12pt]{article}
\usepackage{bmc}
\newcommand{\bandname}{Grades K--1}
\newcommand{\packetid}{NN-K-v1}
\renewcommand{\winrec}{{\large 1st\qquad\quad 2nd}}
\begin{document}
\fontsize{13.5}{18}\selectfont

% ---------------- page 1 ----------------
A break splits one piece into two pieces along a straight line between squares. Break one piece at a time, and never stack pieces.
\gap[3mm]
\begin{center}\breakpicture[7]\end{center}
\gap[6mm]

\prob Break each bar into single squares. Drop a counter in the cup for each break, then write how many counters you used.
\gap[10mm]
\begin{center}
\begin{tabular}{@{}c@{\hspace{16mm}}c@{}}
\chocbar{2}{2}{12} \hspace{6mm} \numbox[16mm] &
\chocbar{2}{3}{12} \hspace{6mm} \numbox[16mm]
\end{tabular}
\\[14mm]
\chocbar{1}{6}{12} \hspace{6mm} \numbox[16mm]
\end{center}

\newpage
% ---------------- page 2 ----------------
\prob Break this bar into single squares using as few breaks as you can. Can anyone at your table do it with fewer breaks?
\gap[6mm]
\begin{center}
\chocbar{3}{4}{12} \hspace{10mm} \numbox[16mm]
\end{center}
\gap[12mm]

\prob Play ``last break wins'' with a partner: take turns making one break, and whoever makes the last break wins. Play two games on each bar, with a different player going first in each game, and circle whether the 1st or 2nd player won.
\gap[8mm]
\begin{center}
\begin{tabular}{@{}c@{\hspace{22mm}}c@{}}
\raisebox{\depth}{\chocbar{3}{3}{12}} & \raisebox{\depth}{\chocbar{2}{4}{12}} \\[9mm]
\begin{tabular}{@{}r@{\quad}l@{}}
\normalsize Game 1 & \winrec\\[5mm]
\normalsize Game 2 & \winrec
\end{tabular}
&
\begin{tabular}{@{}r@{\quad}l@{}}
\normalsize Game 1 & \winrec\\[5mm]
\normalsize Game 2 & \winrec
\end{tabular}
\end{tabular}
\end{center}

\newpage
% ---------------- page 3 ----------------
In 2-Sprigs, each dot starts with two tokens beside it. Players take turns drawing a line from a dot to another dot or back to the same dot. A line may not cross itself or another line, or touch any other dot. Each end of a line uses up one token at its dot, so a loop needs two tokens. Your partner takes away the used tokens. Whoever draws the last line wins.
\gap[4mm]
\begin{center}\sprigsrules\end{center}
\gap[4mm]

\prob Play 2-Sprigs on these boards, and circle whether the 1st or 2nd player won each game.
\gap[4mm]
\begin{center}
\recboard{\board{88}{52}{\Ltwo}}{\winrec}\hfill
\recboard{\board{88}{52}{\Ltwo}}{\winrec}
\\[5mm]
\recboard{\board{88}{76}{\Lthree}}{\winrec}\hfill
\recboard{\board{88}{76}{\Lthree}}{\winrec}
\end{center}

\newpage
% ---------------- page 4 ----------------
\prob Play 2-Sprigs on these boards. When a game is over, count its lines and circle who won.
\gap[5mm]
\begin{center}
\recboard{\board{88}{88}{\Lfour}}{\normalsize lines \numbox[10mm]\qquad \winrec}\hfill
\recboard{\board{88}{88}{\Lfour}}{\normalsize lines \numbox[10mm]\qquad \winrec}
\\[8mm]
\recboard{\board{88}{88}{\Lfour}}{\normalsize lines \numbox[10mm]\qquad \winrec}\hfill
\recboard{\board{88}{88}{\Lfour}}{\normalsize lines \numbox[10mm]\qquad \winrec}
\end{center}

\newpage
% ---------------- page 5 ----------------
\prob Can the 2nd player win 2-Sprigs on a board with 5 dots? Play to find out.
\gap[5mm]
\begin{center}
\recboard{\board{90}{90}{\Lfive}}{\winrec}\hfill
\recboard{\board{90}{90}{\Lfive}}{\winrec}
\\[8mm]
\recboard{\board{90}{90}{\Lfive}}{\winrec}\hfill
\recboard{\board{90}{90}{\Lfive}}{\winrec}
\end{center}

\newpage
% ---------------- page 6 ----------------
\prob Before you play on each board, guess whether the 1st or 2nd player will win. Then play and see.
\gap[5mm]
\begin{center}
\board{140}{82}{\Lsixwide}\\[3mm]
{\normalsize guess}\enspace\winrec\qquad\qquad{\normalsize winner}\enspace\winrec
\\[10mm]
\board{140}{82}{\Lsixwide}\\[3mm]
{\normalsize guess}\enspace\winrec\qquad\qquad{\normalsize winner}\enspace\winrec
\end{center}

\newpage
% ---------------- page 7 ----------------
\prob Each group of pieces was broken from one bar. How many breaks did each group take?
\gap[8mm]
\newcommand{\piece}[4]{% rows cols x y (in squares)
  \begin{scope}[shift={(#3,#4)}]
    \fill[black!10] (0,0) rectangle (#2,#1);
    \draw[step=1,black!65,line width=0.5pt] (0,0) grid (#2,#1);
    \draw[line width=1.1pt] (0,0) rectangle (#2,#1);
  \end{scope}}
\begin{center}
% from a 2 by 3 bar: 2x2, 1x1, 1x1
\begin{tikzpicture}[x=12mm,y=12mm,baseline=(current bounding box.center)]
  \piece{2}{2}{0}{0} \piece{1}{1}{2.6}{0} \piece{1}{1}{4.2}{0}
\end{tikzpicture}
\hspace{16mm}\numbox[16mm]
\\[14mm]
% from a 3 by 3 bar: 3x1, 1x2, 2x1, 1x1, 1x1
\begin{tikzpicture}[x=12mm,y=12mm,baseline=(current bounding box.center)]
  \piece{3}{1}{0}{0} \piece{1}{2}{1.6}{0} \piece{2}{1}{4.2}{0} \piece{1}{1}{5.8}{0} \piece{1}{1}{7.4}{0}
\end{tikzpicture}
\hspace{12mm}\numbox[16mm]
\\[14mm]
% from a 3 by 4 bar: 1x4, 2x1, 1x2, 1x1, 1x1, 1x1, 1x1
\begin{tikzpicture}[x=12mm,y=12mm,baseline=(current bounding box.center)]
  \piece{1}{4}{0}{1.6} \piece{2}{1}{4.6}{1.6} \piece{1}{2}{6.2}{1.6}
  \piece{1}{1}{0.6}{0} \piece{1}{1}{2.2}{0} \piece{1}{1}{3.8}{0} \piece{1}{1}{5.4}{0}
\end{tikzpicture}
\hspace{12mm}\numbox[16mm]
\end{center}

\end{document}
